% =============================================================================
% التوزيع السنوي — المعلوماتية — جذع مشترك علوم وتكنولوجيا
% السنة الدراسية: 2026 / 2027
% =============================================================================
% Compile with XeLaTeX:
%   xelatex -interaction=nonstopmode -halt-on-error -output-directory=build annual-distribution-science.tex
% =============================================================================

\documentclass[10pt, a4paper, landscape]{article}

% ── Geometry ─────────────────────────────────────────────────────────────────
\usepackage[
  top=1.1cm,
  bottom=0.9cm,
  left=1.2cm,
  right=1.2cm
]{geometry}

% ── Packages (fancyhdr & friends BEFORE polyglossia) ─────────────────────────
\usepackage{fancyhdr}
\usepackage{titlesec}
\usepackage{enumitem}
\usepackage{xcolor}
\usepackage{tikz}
\usepackage{eso-pic}
\usepackage{tcolorbox}
\tcbuselibrary{skins, breakable}
\usepackage{array}
\usepackage{longtable}
\usepackage{colortbl}
\usepackage{multirow}
\usepackage{hyperref}

% ── Arabic / RTL ─────────────────────────────────────────────────────────────
\usepackage{polyglossia}
\setmainlanguage[numerals=western]{arabic}
\setotherlanguage{english}

% ── Fonts ────────────────────────────────────────────────────────────────────
\setmainfont{Amiri}[
  BoldFont = * Bold,
  ItalicFont = * Italic,
]
\newfontfamily\arabicfont{Amiri}[
  Script = Arabic,
  BoldFont = * Bold,
]

% ── Colors ───────────────────────────────────────────────────────────────────
\definecolor{primary}{HTML}{1A5276}
\definecolor{secondary}{HTML}{2E86C1}
\definecolor{accent}{HTML}{D4AC0D}
\definecolor{lightbg}{HTML}{EBF5FB}
\definecolor{darktext}{HTML}{1C2833}
\definecolor{rulegray}{HTML}{ABB2B9}
\definecolor{holiday}{HTML}{FDF2E3}

% ── Header / Footer (none) ───────────────────────────────────────────────────
\pagestyle{empty}

% ── Page Border (1cm from all edges) ─────────────────────────────────────────
\AddToShipoutPictureBG{%
  \begin{tikzpicture}[remember picture, overlay]
    \draw[primary, line width=1.2pt]
      ([xshift=-1cm, yshift=-1cm]current page.north east)
      rectangle
      ([xshift=1cm, yshift=1cm]current page.south west);
  \end{tikzpicture}%
}

% ── Table helpers ────────────────────────────────────────────────────────────
\setlength{\arrayrulewidth}{0.4pt}
\setlength{\tabcolsep}{3pt}
\arrayrulecolor{rulegray}
\renewcommand{\arraystretch}{0.9}

\newcommand{\thd}[1]{\cellcolor{primary}\color{white}\textbf{#1}}

% ── Hyperref ─────────────────────────────────────────────────────────────────
\hypersetup{
  colorlinks=false,
  pdfauthor={لفقير عبدالجليل},
  pdftitle={التوزيع السنوي — المعلوماتية — علوم وتكنولوجيا 2026-2027},
}

% ══════════════════════════════════════════════════════════════════════════════
\begin{document}

% ── Title block ──────────────────────────────────────────────────────────────
\begin{center}
  {\scriptsize\color{secondary} الجمهورية الجزائرية الديمقراطية الشعبية \quad -- \quad
   وزارة التربية الوطنية \quad -- \quad مديرية التربية لولاية أدرار \quad -- \quad
   ثانوية الشهيد حكومي العيد أدرار} \\[2pt]
  {\large\bfseries\color{primary} التوزيع السنوي لمادة المعلوماتية — السنة الأولى ثانوي — جذع مشترك علوم وتكنولوجيا}
\end{center}

\vspace{1pt}

% ── Info bar ─────────────────────────────────────────────────────────────────
\begin{tcolorbox}[
  enhanced, colback=lightbg, colframe=secondary, boxrule=0.7pt, arc=3pt,
  left=8pt, right=8pt, top=4pt, bottom=4pt,
]
  \small\centering
  \textbf{الموسم الدراسي:} 2026/2027 \quad|\quad
  \textbf{الأستاذ:} لفقير عبدالجليل \quad|\quad
  \textbf{الحجم الساعي الأسبوعي:} ساعتان \quad|\quad
  \textbf{المعامل:} 02 \quad|\quad
  \textbf{الحجم السنوي:} 48 حصة (24 نظري + 24 تطبيقي)
\end{tcolorbox}

\vspace{3pt}

% ── Distribution table ───────────────────────────────────────────────────────
{\scriptsize
\begin{longtable}{|>{\centering\arraybackslash}p{1.5cm}|>{\centering\arraybackslash}p{0.9cm}|>{\centering\arraybackslash}p{1.2cm}|>{\centering\arraybackslash}p{2.9cm}|>{\centering\arraybackslash}p{3.1cm}|p{3.5cm}|p{6.0cm}|p{6.0cm}|}
\hline
\thd{الأشهر} & \thd{الأسبوع} & \thd{أسبوع العمل} & \thd{التاريخ} & \thd{المجال} & \thd{الدرس النظري} & \thd{موضوع الوحدة} & \thd{العمل التطبيقي} \\
\hline
\endfirsthead
\hline
\thd{الأشهر} & \thd{الأسبوع} & \thd{أسبوع العمل} & \thd{التاريخ} & \thd{المجال} & \thd{الدرس النظري} & \thd{موضوع الوحدة} & \thd{العمل التطبيقي} \\
\hline
\endhead
\hline
\endfoot

\multirow{2}{*}{\parbox{1.3cm}{\centering سبتمبر}} & 1 & 1 & 21 سبتمبر - 24 سبتمبر & \multirow{11}{*}{\parbox{2.9cm}{\centering بيئة التعامل مع الحاسوب}} & — & نظرة عن البرنامج + توجيهات عامة & تقويم تشخيصي \\
\hline
سبتمبر & 2 & 2 & 27 سبتمبر - 01 أكتوبر &  & تقنية المعلومات & مفاهيم حول تقنية المعلومات & تطور تقنية المعلومات \\
\hline
\multirow{4}{*}{\parbox{1.3cm}{\centering أكتوبر}} & 3 & 3 & 04 أكتوبر - 08 أكتوبر &  & تجميع الحاسوب & وحدات الحاسوب & مكونات + التركيب الفعلي \\
\hline
أكتوبر & 4 & 4 & 11 أكتوبر - 15 أكتوبر &  &  & وحدة قياس الذاكرة & محاكاة تركيب الحاسوب \\
\hline
أكتوبر & 5 & 5 & 18 أكتوبر - 22 أكتوبر &  & نظام التشغيل & نظام التشغيل & إعدادات البيوس \\
\hline
أكتوبر & 6 & — & 29 أكتوبر - 08 نوفمبر &  & \multicolumn{3}{c|}{عطلة الخريف} \\
\hline
\multirow{4}{*}{\parbox{1.3cm}{\centering نوفمبر}} & 7 & 6 & 08 نوفمبر - 12 نوفمبر &  & نظام التشغيل & إعدادات التثبيت & محاكاة تثبيت نظام التشغيل \\
\hline
نوفمبر & 8 & 7 & 15 نوفمبر - 19 نوفمبر &  & حماية الحاسوب & حماية الحاسوب & لوحة التحكم \\
\hline
نوفمبر & 9 & 8 & 22 نوفمبر - 26 نوفمبر &  & الشبكة المحلية & مفاهيم عامة حول الشبكات & تركيب وإعداد الشبكة \\
\hline
نوفمبر & 10 & 9 & 29 نوفمبر - 03 ديسمبر &  & الشبكة المحلية & تصنيفات الشبكة & مشاركة موارد باستعمال الشبكة \\
\hline
\multirow{4}{*}{\parbox{1.3cm}{\centering ديسمبر}} & 11 & — & 06 ديسمبر - 10 ديسمبر &  & \multicolumn{3}{c|}{اختبارات الفصل الأول} \\
\hline
ديسمبر & 12 & 10 & 13 ديسمبر - 17 ديسمبر & \multirow{10}{*}{\parbox{2.9cm}{\centering مدخل إلى البرمجة}} & المخطط الانسيابي & مفاهيم حول المخطط الانسيابي & رسم مخطط باستعمال LARP \\
\hline
ديسمبر & 13 & — & 17 ديسمبر - 03 جانفي &  & \multicolumn{3}{c|}{عطلة الشتاء} \\
\hline
ديسمبر & 14 & — & 17 ديسمبر - 03 جانفي &  & \multicolumn{3}{c|}{عطلة الشتاء} \\
\hline
\multirow{5}{*}{\parbox{1.3cm}{\centering جانفي}} & 15 & 11 & 03 جانفي - 07 جانفي &  & المخطط الانسيابي & الشرط & رسم مخطط به شرط \\
\hline
جانفي & 16 & 12 & 10 جانفي - 14 جانفي &  & المخطط الانسيابي & الحلقة & رسم مخطط به حلقة \\
\hline
جانفي & 17 & 13 & 17 جانفي - 21 جانفي &  & مدخل إلى الخوارزميات & هيكل الخوارزمية وقواعدها & استعمال برنامج Algobox \\
\hline
جانفي & 18 & 14 & 24 جانفي - 28 جانفي &  & مدخل إلى الخوارزميات & قواعد المعرفات والإسناد & تطبيقات بها تعليمات الإسناد \\
\hline
جانفي & 19 & 15 & 31 جانفي - 04 فيفري &  & التعليمات الأساسية & التعليمة الشرطية & تطبيقات بها تعليمات شرطية \\
\hline
\multirow{4}{*}{\parbox{1.3cm}{\centering فيفري}} & 20 & 16 & 07 فيفري - 11 فيفري &  & التعليمات الأساسية & التعليمة التكرارية & تطبيقات بها تعليمات تكرارية \\
\hline
فيفري & 21 & 17 & 14 فيفري - 18 فيفري &  & التعليمات الأساسية & حل مسائل مختلفة & حل مسائل مختلفة \\
\hline
فيفري & 22 & 18 & 21 فيفري - 25 فيفري & \multirow{9}{*}{\parbox{2.9cm}{\centering تقنيات الويب}} & المتصفح & المتصفح & استعمال المتصفح \\
\hline
فيفري & 23 & 19 & 28 فيفري - 04 مارس &  & إنشاء صفحة ويب & هيكل صفحة الويب & كتابة هيكل صفحة الويب \\
\hline
\multirow{4}{*}{\parbox{1.3cm}{\centering مارس}} & 24 & — & 07 مارس - 11 مارس &  & \multicolumn{3}{c|}{اختبارات الفصل الثاني} \\
\hline
مارس & 25 & 20 & 14 مارس - 18 مارس &  & إنشاء صفحة ويب & النصوص والعناوين & إنشاء صفحة بها عناوين ونص \\
\hline
مارس & 26 & — & 18 مارس - 04 أفريل &  & \multicolumn{3}{c|}{عطلة الربيع} \\
\hline
مارس & 27 & — & 18 مارس - 04 أفريل &  & \multicolumn{3}{c|}{عطلة الربيع} \\
\hline
\multirow{4}{*}{\parbox{1.3cm}{\centering أفريل}} & 28 & — & 04 أفريل - 08 أفريل &  & إنشاء صفحة ويب & القوائم والصور & إضافة صورة وروابط تشعبية \\
\hline
أفريل & 29 & 22 & 11 أفريل - 15 أفريل &  & إنشاء صفحة ويب & الارتباط التشعبي والجداول & رسم جدول وتنسيقه \\
\hline
أفريل & 30 & 23 & 18 أفريل - 22 أفريل &  & أدوات التواصل & مواقع التواصل الاجتماعي & البريد الإلكتروني \\
\hline
أفريل & 31 & 24 & 25 أفريل - 29 أفريل & \multirow{2}{*}{\parbox{2.9cm}{\centering المكتبية}} & معالج النصوص & الأنماط & تنسيق نصوص باستعمال الأنماط \\
\hline
\multirow{1}{*}{\parbox{1.3cm}{\centering ماي}} & 32 & — & 02 ماي - 06 ماي &  & \multicolumn{3}{c|}{اختبارات الفصل الثالث} \\
\hline
\hline
\end{longtable}
}

\vspace{2pt}

% ── Signatures ───────────────────────────────────────────────────────────────
\vspace{-2pt}
\noindent
\begin{minipage}[t]{0.30\textwidth}
  \centering
  \textbf{\color{darktext} لفقير عبدالجليل الأستاذ} \\[0.55cm]
  {\color{rulegray}\hrulefill} \\[2pt]
  {\small لفقير عبدالجليل}
\end{minipage}%
\hfill
\begin{minipage}[t]{0.30\textwidth}
  \centering
  \textbf{\color{darktext}المدير} \\[0.55cm]
  {\color{rulegray}\hrulefill}
\end{minipage}%
\hfill
\begin{minipage}[t]{0.30\textwidth}
  \centering
  \textbf{\color{darktext}المفتش} \\[0.55cm]
  {\color{rulegray}\hrulefill}
\end{minipage}

\end{document}
